\documentclass[11pt]{article}
\usepackage{template}
\usepackage{amsmath,amssymb,microtype,tikz,graphicx,xurl}
\usetikzlibrary{arrows.meta,positioning,calc}
\setlength{\emergencystretch}{2.5em}
\raggedbottom
\SetKeys{date=2026-10-27,vol=014,title={每日有趣新闻｜时政、科学与数学},author={ChatGPT 编辑部},note={\small 采编窗口：2026年10月25—27日；截至北京时间10月27日上午。日期分别标明事件发生与报道发布；未经确认的政治指控、预印本结论及市场预期均不作为既成事实。}}
\renewcommand{\newsitem}[3]{\par\medskip{\large\bfseries\color{newsblue}#1\par}{\footnotesize\color{newsgray}#2\par}\smallskip\color{black}#3\par\medskip{\color{newsgray}\rule{\linewidth}{0.25pt}}\par}
\newcommand{\newslead}[3]{\par\medskip{\LARGE\bfseries\color{newsblue}#1\par}\vspace{1mm}{\small\color{newsgray}#2\par}\smallskip\color{black}#3\par\medskip{\color{newsgray}\rule{\linewidth}{0.45pt}}\par}
\newcommand{\newssub}[1]{\par\smallskip\noindent{\bfseries\color{newsblue}#1}\quad}
\begin{document}
\begin{documents}
\raggedcolumns

\newslead{Anthropic公开招股书：营收十四倍增长，807亿美元净亏损究竟意味着什么？}{AI产业与资本市场｜10月26日美国SEC受理S-1注册文件；数据截至2026年6月30日}{
10月26日，开发Claude模型的Anthropic向美国证券交易委员会（SEC）公开提交S-1注册声明，拟在纳斯达克全球精选市场以“ANT”为代码发行A类普通股。SEC的EDGAR档案显示，申请主体为Anthropic, PBC，申报编号333-299633，文件当日受理。\textbf{提交招股书不等于完成上市}：发行股数、定价区间与最终募集金额尚未确定，所谓1.5万亿至2万亿美元估值、可能募集1000亿美元，仍是媒体援引市场人士的讨论，并非已获确认的发行结果。\cite{AnthropicSEC,AnthropicAP}

\newssub{先拆开两组容易误读的数字。}招股书披露，2026年上半年营业收入约162.86亿美元，而2025年同期约11.63亿美元，增幅约十四倍；同期报告的净亏损却达到807.48亿美元。单看这两个数字，似乎企业每赚一美元就亏五美元，但会计结构远比这种比较复杂。约723.88亿美元的损失来自可转换融资工具及相关负债公允价值变化，这类非现金会计费用与当期向员工支付工资、购买算力产生的经营成本并不等价。公司同期营业亏损约6.93亿美元，明显小于报告净亏损。\cite{AnthropicSEC,AnthropicFinance,AnthropicAP}

\begin{center}
\begin{tikzpicture}[x=0.31cm,y=0.34cm,font=\scriptsize]
\draw[gray!50] (0,0)--(19,0);
\node[anchor=west,font=\bfseries] at (0,8.7) {Anthropic 2026上半年财务数字（亿美元）};
\node[anchor=west] at (0,7.4) {营业收入};
\fill[newsblue!25] (5.2,6.7) rectangle (5.9,7.4);
\node[anchor=west] at (6.2,7.05) {2025同期：11.63};
\fill[newsblue!75] (5.2,5.7) rectangle (14.9,6.4);
\node[anchor=west] at (15.1,6.05) {2026：162.86};
\node[anchor=west] at (0,4.55) {报告净亏损};
\fill[red!60] (5.2,3.75) rectangle (17.7,4.45);
\node[anchor=west] at (18,4.1) {807.48};
\node[anchor=west] at (0,2.75) {其中非现金公允价值变动};
\fill[orange!65] (5.2,1.9) rectangle (16.4,2.6);
\node[anchor=west] at (16.7,2.25) {723.88};
\node[anchor=west,text=gray!80!black] at (0,0.7) {长度仅在同一行组内示意，跨组不可按长度比较};
\end{tikzpicture}
\par\footnotesize\color{newsgray}图1　编辑部原创示意。金额来自2026年10月26日SEC S-1；收入与净亏损属不同会计项目，图形非统一比例坐标。\end{center}\color{black}

这也不意味着公司已经能够不依赖外部融资而持续扩张。其计算与基础设施成本上半年约89.45亿美元，约占同期收入55\%；云计算、芯片、数据中心和长期容量承诺构成了未来成本的重要来源。收入增长快，说明客户付费规模在扩大；经营亏损收窄，说明当前成本与收入的关系有所改善；巨额长期承诺则说明即使增长速度放缓，未来仍可能面临较高的固定支付压力。三个命题必须分别检验。\cite{AnthropicSEC,AnthropicFinance}

\newssub{上市究竟改变了什么？}首次公开发行能够扩大股权融资渠道，并让更多投资者获得定期披露的信息，但不自动降低模型训练和推理的单位成本。投资者真正需要观察的是经常性收入能否留存、客户是否过于集中、模型价格下降是否快于计算成本改善，以及巨额合同承诺是否有相匹配的未来需求。公司在文件中提示其五大客户贡献了约四分之一收入，说明少数大型企业采购对营收有显著影响。与此同时，模型厂商可能同时是云服务商的客户、合作伙伴与竞争者，不能把产业链中的所有合同金额简单相加当作新增终端需求。\cite{AnthropicSEC,AnthropicGuardian}

\newssub{治理结构与安全承诺。}Anthropic拟通过不同类别股份让七位联合创始人掌握约50.1\%投票权，另设长期公益信托参与董事任命。这样的结构可能有助于维持创始团队的长期研发方向，却也意味着普通A类股东的表决影响力与经济权益不完全对称。文件还将模型滥用、网络安全、法律责任、监管以及极端AI事故列为风险；这属于上市文件要求的风险披露，不能因为写进招股书就认为每种情景已经发生，也不能因为公司声称重视安全就认为风险已经被解决。\cite{AnthropicSEC,AnthropicAP}

对这家以安全研究作为品牌特色的企业而言，公开市场带来一个值得持续观察的张力：更高的透明度、融资能力与增长要求，将如何与模型安全评估、产品限制和社会责任相协调？这个问题没有在10月26日得到答案。下一步应阅读正式定价文件、季度财务更新和实际治理安排，而不是把拟议估值当作公司的已实现价值。
}

\newslead{埃博拉疫情跨境传播：近万确诊背后，追踪接触者为何比边境封锁更关键？}{公共卫生深度｜世卫组织10月23日疫情通报，统计截至10月20日；10月26日新增转运报道}{
世界卫生组织（WHO）10月23日通报，刚果民主共和国自5月以来的邦迪布焦型埃博拉病毒病疫情持续扩散。截至10月20日，该国累计报告\textbf{9704例确诊、4687例死亡}，粗病死比例为48.3\%，涉及八个省的66个卫生区。相较10月6日的8728例确诊、4205例死亡，14天内累计数字分别增加976例与482例。这里的“粗病死比例”是报告死亡数除以报告确诊数，受到报告延迟、病例发现与治疗可及性的影响，不等于每位感染者面临完全相同的死亡概率。\cite{WHOOct23,WHOOct8}

\begin{center}
\begin{tikzpicture}[x=0.36cm,y=0.46cm,font=\scriptsize]
\node[anchor=west,font=\bfseries] at (0,6.6) {刚果（金）邦迪布焦型疫情：两次通报累计数};
\node[anchor=west] at (0,5.55) {10月6日确诊};
\fill[newsblue!45] (5,5.1) rectangle (13.73,5.7);
\node[anchor=west] at (14,5.4) {8728};
\node[anchor=west] at (0,4.45) {10月20日确诊};
\fill[newsblue!85] (5,4.0) rectangle (14.70,4.6);
\node[anchor=west] at (15,4.3) {9704};
\node[anchor=west] at (0,3.2) {10月6日死亡};
\fill[red!40] (5,2.75) rectangle (9.21,3.35);
\node[anchor=west] at (9.5,3.05) {4205};
\node[anchor=west] at (0,2.1) {10月20日死亡};
\fill[red!75] (5,1.65) rectangle (9.69,2.25);
\node[anchor=west] at (10,1.95) {4687};
\node[anchor=west,text=gray!80!black] at (0,0.5) {条长以1000例/单位表示；仅统计刚果（金）确诊病例};
\end{tikzpicture}
\par\footnotesize\color{newsgray}图2　编辑部据WHO 10月8日及23日通报绘制。日期为病例统计截止日，非通报发布日期；新增数字不可视为全部发生在14天内。\end{center}\color{black}

\newssub{疫情跨境，不等于各国出现持续社区传播。}WHO确认，肯尼亚10月6日报告一例由刚果经乌干达输入的病例，至10月20日累计报告九例；尼日利亚10月20日也确认一例有乌干达旅行史的输入病例。10月26日，美联社报道尼日利亚患者经医疗转运抵达瑞士，瑞士卫生机构安排隔离与治疗，尼日利亚继续监测相关接触者。肯尼亚和尼日利亚的输入病例表明跨境交通可以带来传播风险，却不能据此断言这些国家都已发生大规模社区流行。\cite{WHOOct23,APEvac}

\newssub{为什么是“21天”？}埃博拉病毒病潜伏期通常为2至21天。对有明确暴露史者开展21天症状监测，是尽早发现病例、减少家庭和医院二次传播的重要手段。监测是否覆盖足够多的接触者、是否能迅速转诊疑似患者，以及医护人员能否获得个人防护装备，往往比笼统宣布关闭边境更能说明防疫体系是否有效。WHO此次明确不建议对受影响国家普遍实施旅行或贸易限制，主张加强口岸监测、实验室诊断、跨境通报与接触者追踪。\cite{WHOOct23}

疫情数字还反映了医疗资源与社会信任问题。WHO报告，最近一轮新增死亡中有相当一部分发生在社区而非治疗中心，意味着不少患者未能及时获得专业照护。10月26日，美联社还报道刚果东部一名参与埃博拉应对的护士遭袭身亡，另一名医护人员受重伤。当地官员尚未确认袭击者身份和具体动机，因此不能把每一次暴力事件都归因于同一组织；但不信任、错误信息、交通困难与长期武装冲突确实使卫生队伍难以进入部分地区。\cite{WHOOct23,APHealth}

\newssub{关于疫苗与风险沟通。}邦迪布焦病毒与2014—2016年西非疫情广为人知的扎伊尔型埃博拉病毒属于不同病毒种。针对一种病毒有效的疫苗，不可未经试验就推定对另一种同样有效。世卫组织此前的风险评估指出，针对本次病毒的获许可疫苗仍然缺乏，这使早发现、隔离、支持性治疗、感染控制和社区参与更加重要。\cite{WHORisk}

这场疫情既是感染病学问题，也是基础公共服务能否运转的问题。新闻报道应避免用“近半数确诊者死亡”制造无差别恐慌：统计数字确实严重，但不同地区、就医时机与病例发现方式会影响观察到的结果。未来更值得追踪的指标包括新增病例趋势、社区死亡占比、接触者随访覆盖率、卫生机构安全及跨境传播链是否被及时截断。
}

\newsitem{巴西大选后股汇齐跌：市场反应不等于选举结果的政治评价}{拉美政治经济｜10月25日总统决选；10月26日路透社市场报道}{
巴西总统卢拉在10月25日决选中获得连任。10月26日，巴西Bovespa股指收跌13.9\%，创下2020年3月以来最大单日百分比跌幅；巴西雷亚尔兑美元收跌约8.1\%，至约5.415雷亚尔兑一美元，盘中跌幅曾更大。巴西交易所B3在股指跌幅触及10\%后暂停交易30分钟；国债市场也承受抛售压力，财政部取消原定部分债券拍卖。\cite{ReutersBrazil}

这些数字首先说明的是\textbf{资产价格对预期变化的快速反应}。第一轮投票后，部分投资者押注右翼候选人胜出并采取更紧缩的财政政策，相关预期此前已反映在股票、汇率与利率期货中。当选举结果与这一押注不一致时，资产价格回调并不意外。但交易市场代表的是有资金参与定价的投资者，而不是全体选民的政治判断；股价跌幅不能直接换算成政策优劣，也不代表巴西经济产出当日实际减少13.9\%。\cite{ReutersBrazil}

下一届政府需要处理公共债务、社会项目、就业与投资之间的财政约束。卢拉团队提出到2030年实现相当于GDP约1.3\%的预算盈余目标，但其具体政策路径仍待公布。市场关心财政部长人选、支出控制与税收安排；普通家庭则更直接关心就业、实际工资、物价和公共服务。不同群体对同一预算方案可能有不同利益与评价，报道不应把任何一方的目标当成唯一尺度。\cite{ReutersBrazil}

另一个未决因素是败选方对结果提出的挑战。政治争议应由巴西选举制度及司法程序处理，而不是由市场波动代替法律认证。后续应观察正式计票认证、行政交接和预算文件，而非只盯住选举次日的价格。
}

\newsitem{英国制裁苏丹冲突双方15名指挥人员：惩责与救援如何并行？}{非洲与人权｜英国政府10月26日公告；路透社交叉核对}{
英国外交部门10月26日宣布，对苏丹武装部队（SAF）、快速支援部队（RSF）及相关武装团体中的15名指挥官和安全人员实施制裁，指控他们涉及性暴力、酷刑和针对平民的袭击。制裁名单是英国政府的法律与外交行动，\textbf{不等于每一名被列入者都已由独立法院定罪}。英国同时宣布追加1500万英镑，用于人道救援和幸存者支持。\cite{UKSudan,ReutersSudan}

事件发生在北达尔富尔首府法希尔陷落一周年之际。该城此前经历长期围困，出现饥饿、人口流离失所及严重侵害平民的报告。英国援引受其资助的调查机构研究，认为交战双方都涉及有系统的剥夺食物与冲突相关性暴力。政府资助的调查有可核查的材料价值，但其结论仍应结合联合国、人道组织与其他独立调查交叉评估；对冲突各方也不应因其政治身份不同而采用不同的证据标准。\cite{UKSudan,ReutersSudan}

制裁通常通过资产冻结、旅行限制或金融交易约束提高涉事人员的成本，不能直接保证战区物资抵达平民手中。若医疗设施受损、运输通道不安全或救援组织无法进入，即使拨款增加，援助也可能无法转化为实际服务。因此评估政策效果至少需要两套指标：一是制裁执行、资金与网络限制；二是物资进入、伤者救治、幸存者保护和调查取证是否改善。\cite{UKSudan}
}

\newsitem{中国召开五中全会：党内治理议程与经济预期需要分开阅读}{中国政治｜10月26日会议开幕；新华社表述经路透社转述核对}{
中国共产党中央委员会10月26日在北京召开为期四天的第五次全体会议。新华社发布的会议材料称，习近平向会议说明关于全面从严治党、加强制度执行的文件草案，强调完善规章制度、监督权力运行与持续反腐。会议尚在进行，最终通过的文件和具体人事决定应以后续正式公报为准。\cite{ReutersPlenum}

这一会议容易被外界与2027年领导层换届、经济政策及反腐整顿联系起来，但这些联系的性质不同：文件草案与会议时间属于可核实事实；政治人事安排的可能变化属于观察者分析；经济增长和资本市场的具体反应则需要数据检验。路透社引用伦敦国王学院学者的看法，认为会议的政治形象展示可能占重要位置。该评论是具名分析，而非中央委员会正式说明。\cite{ReutersPlenum}

官方把制度执行和反腐视为提升治理能力、改善可预期性的途径；另一些观察者关注高层整肃对行政连续性、军队组织和官员激励的影响。这些问题可以并列考察，但不能预设任何一种政治动机已被证明。会议闭幕后，真正可检验的材料是公报文本、制度修订、人事决定和后续执行结果。
}

\newsitem{数学拾趣｜给图染色，为什么还要记录究竟用了几种颜色？}{组合数学｜arXiv:2610.31785；10月26日新列表收录，原始提交时间10月22日；未经同行评议}{
给一张图的顶点染色，要求每条边两端颜色不同，这叫\textbf{正常染色}。例如路径$P_4$有四个依次连接的顶点，用两种颜色交替涂色就满足要求；中心连接三个叶子的星形图$K_{1,3}$同样能用两种颜色正常染色。两张图都“二色可染”，却有完全不同的局部结构。仅仅问最少需要几种颜色，显然不足以识别图的形状。\cite{ArxivColor}

\begin{center}
\begin{tikzpicture}[scale=0.75,every node/.style={circle,draw=black!50,minimum size=0.48cm,inner sep=0pt},line width=0.6pt]
\node[fill=newsblue!65] (a) at (0,0) {};
\node[fill=orange!65] (b) at (1.4,0) {};
\node[fill=newsblue!65] (c) at (2.8,0) {};
\node[fill=orange!65] (d) at (4.2,0) {};
\draw (a)--(b)--(c)--(d);
\node[draw=none,font=\scriptsize] at (2.1,-0.7) {$P_4$：路径};
\node[fill=newsblue!65] (e) at (7.1,0) {};
\node[fill=orange!65] (f) at (7.1,1.2) {};
\node[fill=orange!65] (g) at (5.8,-0.75) {};
\node[fill=orange!65] (h) at (8.4,-0.75) {};
\draw (e)--(f) (e)--(g) (e)--(h);
\node[draw=none,font=\scriptsize] at (7.1,-1.35) {$K_{1,3}$：星形图};
\end{tikzpicture}
\par\footnotesize\color{newsgray}图3　编辑部原创二色正常染色示例。两图都能用两色，但连接结构不同；本图不是论文计算实验的原始图。\end{center}\color{black}

图的色多项式会统计使用给定颜色集合时有多少种正常染色。Stanley提出的\textbf{色对称函数}更进一步：它用对称函数编码各种颜色分配，而不仅仅记录染色总数。数学家一直研究这类不变量是否足以区分某些图，尤其是树。Bingham等六位作者在新预印本中引入带参数$t$的变体$Z_G(\mathbf{x};t)$，用$t$记录实际使用的不同颜色数量，并研究其\textbf{主特殊化}——把变量按照规定取值后得到的较易计算的表达式。\cite{ArxivColor}

作者强调一个微妙区别：他们的$t$-色对称函数作为完整图不变量，与已有色对称函数在信息量上等价；真正增加区分能力的是\textbf{某些特殊化之后仍然保留了更多结构信息}。论文证明，第三主特殊化可以区分并重构一类脊柱只有三个顶点的毛毛虫树，以及某些加长双星图；第四主特殊化能否区分所有毛毛虫树，作者提出了基于计算证据的猜想。\cite{ArxivColor}

这不是“已经证明所有树都可以通过色对称函数唯一识别”，也不是新参数凭空创造了超出原函数的全部信息。它展示的是研究数学不变量的一种策略：原对象可能很复杂，适当选择代数表达与代入方式，可以让原本隐藏的结构在可计算的投影中显现。本刊只核查了论文摘要和定理范围，没有独立检查全部证明，因此以上成果均以作者预印本的证明主张呈现。\cite{ArxivColor}
}

\newsitem{从太空追踪秋色：PACE卫星观测到的“红”并不全是枫叶}{地球科学与文化｜NASA Earth Observatory 10月26日发布}{
美国NASA地球观测站10月26日发布利用PACE卫星追踪秋季植被变色的专题。其OCI高光谱仪在286个光谱波段观测地表和海洋，通过对叶绿素、类胡萝卜素、花青素相关光谱信号的分析，研究植被从绿转红、转黄的季节变化。9月29日至10月4日的全球图像中，高纬度地区已有较强红色信号；10月5日至11日的观测显示，北美中纬度和亚洲部分地区的变色范围发生变化。\cite{NASAFoliage}

这种“红度”是根据遥感光谱构造的指标，不是直接数出每一片红叶。农作物成熟、植被干旱胁迫和云层遮挡也可能影响结果，因此NASA同时邀请公众提交地面照片，帮助校验卫星信号。相比只拍摄一张漂亮的秋景，重复、定量地观察色素变化有助于研究植物物候、碳循环与环境胁迫。\cite{NASAFoliage}

这项研究与东亚的“红叶狩”文化形成有趣对照：古人依靠经验寻找最合适的赏叶时节，现代研究者试图建立全球范围的时间序列。二者目的不同，却都围绕同一件事：秋色并非静态景观，而是一段短暂、随纬度与天气移动的过程。NASA网页提供了可公开查阅的原始卫星图与图注；本期因未取得适于离线归档的原始高分辨率素材，没有用低清缩略图冒充正式科学插图。\cite{NASAFoliage}
}

\newsitem{乌克兰南部局部停火面临落实考验，基辅再遭袭击}{国际简讯｜路透社10月26日，交战方指控未经独立核实}{
俄乌双方此前就赫尔松州俄方控制的奥列什基镇达成局部人道停火安排，计划10月27日开始实施，以便运送食品药品、撤离平民。俄罗斯国防部10月26日指控乌军近期在该镇附近发动炮击与无人机袭击，路透社明确表示无法独立核实这一数字，乌方当时也未就该次指控作出回应。局部停火并不代表整个战场停战，是否真正建立安全走廊还须观察执行。与此同时，路透社报道10月26日晚基辅遭俄军导弹袭击，乌方官员称至少一人死亡、十人受伤，救援仍在进行。\cite{ReutersOleshky,ReutersKyiv}
}

\end{documents}
\vspace{1mm}
{\small\bfseries\color{newsblue}参考资料与编辑说明}\par
{\footnotesize 截稿时间：2026年10月27日北京时间上午。政经材料以政府、监管文件和可信媒体为基础；WHO疫情统计注明截止日；数学为尚未经本刊独立核验的arXiv预印本。三张配图均由编辑部用TikZ原创绘制，不使用合成新闻现场照片。}\par
\begingroup\footnotesize
\renewcommand{\refname}{参考文献}\bibliographystyle{plainurl}
\bibliography{refs}
\endgroup
\end{document}